% /root/monodromy, 2026-11-19. Candidate until parent audit.
\subsection{An explicit invariant covector in every even supersingular degree}
\label{subsec:ss-even-covector}
\status{Candidate specialization of the proved adjunction-square construction;
the complete argument below removes all asymptotic regularity conditions
for this fixed surface.}

\begin{proposition}\label{prop:ss-even-degree-covector}
Let $X=E_0\times E_0$, $E_0:y^2+y=x^3$, and
$L=\mathcal O_{E_0}(3O)\boxtimes\mathcal O_{E_0}(3O)$, over any perfect
field of characteristic two. For every even $d\ge2$, the relative
Hasse--Witt local system on its maximal-rank smooth open has a nonzero
covector defined over the base field. Its rank is $r=9d^2-3$ by
Theorem~\ref{thm:DE-A3-ss-product-rank}, and both images are contained in
\[
 \left\{\begin{pmatrix}A&b\\0&1\end{pmatrix}:
 A\in\GL_{r-1}(\mathbf F_2),\ b\in\mathbf F_2^{r-1}\right\}.
\]
In particular the representation is reducible in every even degree.
This proposition gives an upper bound for the full images, not an equality.
\end{proposition}

\begin{proof}
Put $m=3d$. The line bundle
$N_m=\mathcal O(mO)\boxtimes\mathcal O(mO)$ is the Frobenius pullback
of $N_{m/2}$ on the Frobenius twist. Its canonical Cartier connection
therefore lifts the two invariant vector fields $\delta_1,\delta_2$
to commuting endomorphisms of $H^0(X,N_m)$: in local pullback frames
the transition functions are squares and are killed by these derivations.
On the affine elliptic chart, $\delta x=1$ and $\delta y=x^2$,
so $\delta^2=0$. Consequently
\[
          D=\delta_1\delta_2\quad\text{satisfies }D^2=0.
\]
The invariant top differential is $\omega=dx_1\wedge dx_2$.

In the rational function field, $x_1,x_2$ form a $2$-basis.
Writing a section in a pullback frame as
\[
 s=a^2+x_1b^2+x_2c^2+x_1x_2e^2
\]
gives $D(s)=e^2$ and $C(s\omega)=e\omega$ in that frame.
The Cartier projection formula gives
\begin{equation}\label{eq:ss-even-fixed-Cartier}
              C(sD(s)\omega)=D(s)\omega.
\end{equation}
The local identities glue by the pullback-frame construction.
This is exactly the fixed-covector identity of
Theorem~\ref{thm:general-adjunction-covector}: under Serre duality,
$D(s)\omega$ gives a regular Frobenius-equivariant covector on the
full coherent bundle, whose restriction to its relative kernel kills
the nilpotent part and takes values in $\mathbf F_2$ on its Lang
local system. The dependence on $s$ is linear, so the construction
descends from the coefficient cone to the projective parameter space.
Equivalently, on a smooth section it is the differential
$\operatorname{Res}_Y(D(s)\omega/s)$; the residue is unchanged by
scaling $s$.

This covector is nonzero generically without a regularity hypothesis.
The section $x_1x_2\in H^0(X,N_m)$ is defined for $m\ge6$ and
$D(x_1x_2)=1$, so $D$ is a nonzero endomorphism. For generic $s$,
$D(s)\ne0$. It cannot then be a scalar multiple of $s$: if
$D(s)=c s$, applying $D$ gives $0=c^2s$, hence $c=0$.
The kernel of the adjunction map $H^0(N_m)\omega\to H^0(\omega_Y)$
is precisely $k s\omega$, because $H^0(\omega_X)=k\omega$.
Thus the relative covector is nonzero generically. Its rank as a
morphism of finite local systems is locally constant on the connected
maximal-rank open, so it is nonzero everywhere there. The entire
construction is defined over $\mathbf F_2$, and hence over the stated
base field. The displayed pointwise-stabilizer containment follows.
\end{proof}
